\begin{ouverture}[suites]{10}{Puissances et suites matricielles}{Graphes et matrices}{Récurrence · Point fixe · Ordres 2 et 3}
\objectif{Calculer des puissances de matrices en ordres $2$ et $3$.}
\objectif{Démontrer une formule par récurrence.}
\objectif{Modéliser une évolution couplée par colonnes.}
\objectif{Résoudre une récurrence affine avec un point fixe.}
\prerequis{Calcul matriciel, suites et récurrence.}
\end{ouverture}
\essentiel

\essSection{1}{Répéter une transformation}
Pour une matrice carrée, $A^{n+1}=A^nA$ et $A^0=I$. Pour une matrice colonne $U_{n+1}=AU_n$, une récurrence donne $U_n=A^nU_0$.
On ne calcule pas $A^n$ en élevant chaque coefficient à la puissance $n$.
Pour $D=\operatorname{diag}(a,b)$, en revanche $D^n=\operatorname{diag}(a^n,b^n)$.
\begin{exemple}[Une matrice de cisaillement]
Pour $A=\begin{pmatrix}1&1\\0&1\end{pmatrix}$,
$A^n=\begin{pmatrix}1&n\\0&1\end{pmatrix}$.
Au rang zéro, on retrouve $I$ ; multiplier la formule au rang $n$ par $A$ remplace le coefficient $n$ par $n+1$.
\end{exemple}
\essSection{2}{Conjecturer puis prouver en ordre $2$}
Pour $B=\begin{pmatrix}3&1\\0&2\end{pmatrix}$, calculer $B^2,B^3$ suggère
$B^n=\begin{pmatrix}3^n&3^n-2^n\\0&2^n\end{pmatrix}$.
Le coefficient supérieur droit du produit au rang suivant vaut $3^n+2(3^n-2^n)=3^{n+1}-2^{n+1}$, ce qui établit l'hérédité.
Une autre méthode guidée consiste à fournir $B=PDP^{-1}$. Alors $(PDP^{-1})^n=PD^nP^{-1}$, puisque les facteurs intermédiaires $P^{-1}P$ se simplifient. Le manuel ne suppose pas une théorie générale de diagonalisation.
\essSection{3}{Un exemple en ordre $3$}
Pour $T=\begin{pmatrix}1&1&0\\0&1&1\\0&0&1\end{pmatrix}$, on a
\[T^n=\begin{pmatrix}1&n&n(n-1)/2\\0&1&n\\0&0&1\end{pmatrix}.
\]
Pour vérifier l'hérédité, le coefficient $(1,3)$ devient $n(n-1)/2+n=n(n+1)/2$. Les autres coefficients se contrôlent directement.
Cette formule illustre la présence naturelle de suites arithmétiques et de coefficients binomiaux dans les puissances de matrices.
\essSection{4}{Évolutions affines couplées}
\begin{savoirfaire}
Pour $U_{n+1}=AU_n+C$, chercher une colonne constante $L$ vérifiant $(I-A)L=C$.
Si une telle colonne existe, $V_n=U_n-L$ satisfait $V_{n+1}=AV_n$, donc
$U_n=L+A^n(U_0-L)$.
Si $I-A$ est inversible, $L=(I-A)^{-1}C$ est unique.
\end{savoirfaire}
\begin{exemple}
Pour $A=\begin{pmatrix}1/2&0\\0&1/4\end{pmatrix}$, $C=\binom13$ et $U_0=\binom00$, on trouve $L=\binom24$.
Ainsi $U_n=\binom{2(1-2^{-n})}{4(1-4^{-n})}$, qui tend vers $L$ coordonnée par coordonnée.
\end{exemple}
\figurelivre{e10-evolution.pdf}{Les deux composantes de l'exemple tendent vers leurs valeurs d'équilibre, $2$ et $4$. Le modèle matriciel regroupe les deux récurrences et conserve leur interprétation.}
\essSection{5}{La matrice ne garantit pas une limite}
Une suite matricielle peut diverger ou être périodique. Pour $A=\begin{pmatrix}0&1\\1&0\end{pmatrix}$, $A^2=I$ : certaines colonnes alternent indéfiniment.
Des modèles de populations peuvent coupler plusieurs effectifs. Il faut vérifier la pertinence des coefficients et le maintien de valeurs admissibles ; une formule algébrique n'assure pas qu'une approximation physique reste valable sur une longue durée.
Le programme distingue ces colonnes d'évolution des distributions lignes des chaînes de Markov, étudiées au chapitre 12.

\exercices
\markright{Prendre confiance}
\begingroup\renewcommand{\otBesoinPlace}[1]{}\begin{minipage}{\linewidth}
\rubrique{Prendre confiance}
\begin{exercice}[1]{Puissances diagonales}
$D=\begin{pmatrix}2&0\\0&3\end{pmatrix}$.
\begin{questions}
\item Calculer $D^2$.
\item Donner $D^n$.
\item Donner $D^{-1}$.
\end{questions}
\end{exercice}
\begin{exercice}[1]{Matrice périodique}
$A=\begin{pmatrix}0&1\\1&0\end{pmatrix}$.
\begin{questions}
\item Calculer $A^2$.
\item Donner $A^{2026}$.
\item Donner $A^{2027}$.
\end{questions}
\end{exercice}
\begin{exercice}[1]{Premières colonnes}
$A=\begin{pmatrix}1&1\\0&1\end{pmatrix}$ et $U_0=\binom01$.
\begin{questions}
\item Calculer $U_1=AU_0$.
\item Calculer $U_2$.
\item Conjecturer $U_n$.
\end{questions}
\end{exercice}
\begin{exercice}[1]{Équilibre}
$U_{n+1}=\frac12I_2U_n+\binom12$.
\begin{questions}
\item Chercher la colonne fixe $L$.
\item Exprimer $U_n-L$.
\item Si $U_0=0$, donner la limite.
\end{questions}
\end{exercice}
\end{minipage}\endgroup
\clearpage
\markright{Appliquer et relier}
\begingroup\renewcommand{\otBesoinPlace}[1]{}\begin{minipage}{\linewidth}
\rubrique{Appliquer et relier}
\begin{exercice}[2]{Prouver une puissance}
$A=\begin{pmatrix}1&1\\0&1\end{pmatrix}$.
\begin{questions}
\item Verifier la formule du cours pour $n=0$.
\item Établir l’hérédité par multiplication.
\item Donner $A^{100}$.
\end{questions}
\end{exercice}
\begin{exercice}[2]{Une puissance triangulaire}
$B=\begin{pmatrix}3&1\\0&2\end{pmatrix}$.
\begin{questions}
\item Calculer $B^2$.
\item Donner $B^n$ et vérifier au rang $1$.
\item Calculer $B^5$.
\end{questions}
\end{exercice}
\begin{exercice}[2]{Ordre $3$}
$T=\begin{pmatrix}1&1&0\\0&1&1\\0&0&1\end{pmatrix}$.
\begin{questions}
\item Calculer $T^2$.
\item Utiliser la formule du cours pour $T^4$.
\item Calculer $T^4\begin{pmatrix}0\\0\\1\end{pmatrix}$.
\end{questions}
\end{exercice}
\begin{exercice}[2]{Deux récurrences}
$u_{n+1}=u_n/2+1$, $v_{n+1}=v_n/4+3$, $u_0=v_0=0$.
\begin{questions}
\item Écrire $U_{n+1}=AU_n+C$.
\item Trouver $L$.
\item Exprimer les deux composantes et leurs limites.
\end{questions}
\end{exercice}
\end{minipage}\endgroup
\clearpage
\markright{Chercher et argumenter}
\begingroup\renewcommand{\otBesoinPlace}[1]{}\begin{minipage}{\linewidth}
\rubrique{Chercher et argumenter}
\begin{exercice}[3]{Une méthode avec $PDP^{-1}$}
$P=\begin{pmatrix}1&-1\\0&1\end{pmatrix}$ et $D=\operatorname{diag}(3,2)$.
\begin{questions}
\item Calculer $P^{-1}$ et $PDP^{-1}$.
\item Justifier $(PDP^{-1})^n=PD^nP^{-1}$.
\item Retrouver les puissances de $B$.
\end{questions}
\end{exercice}
\begin{exercice}[3]{Un couplage périodique}
$A=\begin{pmatrix}0&2\\1/2&0\end{pmatrix}$, $U_0=\binom10$.
\begin{questions}
\item Calculer $A^2$.
\item Donner les colonnes paires et impaires.
\item La suite converge-t-elle ?
\end{questions}
\end{exercice}
\begin{exercice}[3]{Point fixe non unique}
$A=I_2$ et $C=0$.
\begin{questions}
\item Quelles colonnes sont fixes ?
\item Que vaut $U_n$ ?
\item Pourquoi ne peut-on utiliser $(I-A)^{-1}C$ ?
\end{questions}
\end{exercice}
\begin{exercice}[3]{Contrôler une formule en ordre $3$}
On propose la formule de $T^n$ du cours.
\begin{questions}
\item Verifier les coefficients pour $n=0,1$.
\item Calculer le coefficient $(1,3)$ de $T^nT$.
\item Expliquer l’apparition de $\binom n2$.
\end{questions}
\end{exercice}
\end{minipage}\endgroup
